\def\stat{konashenkova}





\def\tit{МОДЕЛИРОВАНИЕ НЕСТАЦИОНАРНОГО СТОХАСТИЧЕСКОГО ПРОЦЕССА\\ ПОСРЕДСТВОМ 


ЕГО~КАНОНИЧЕСКОГО РАЗЛОЖЕНИЯ\\ НА~ОСНОВЕ ВЕЙВЛЕТ-НЕЙРОННОЙ СЕТИ$^*$}





\def\titkol{Моделирование нестационарного СтП %стохастического процесса 


посредством  его~КР на~основе ВНС} %вейвлет-нейронной сети}





\def\aut{И.\,Н.~Синицын$^1$, В.\,И.~Синицын$^2$, Э.\,Р.~Корепанов$^3$, 


Т.\,Д.~Конашенкова$^4$}





\def\autkol{И.\,Н.~Синицын, В.\,И.~Синицын, Э.\,Р.~Корепанов, 


Т.\,Д.~Конашенкова}





\titel{\tit}{\aut}{\autkol}{\titkol}





\index{Синицын И.\,Н.}


\index{Синицын В.\,И.}


\index{Корепанов Э.\,Р.}


\index{Конашенкова Т.\,Д.}


\index{Sinitsyn I.\,N.}


\index{Sinitsyn V.\,I.}


\index{Korepanov E.\,R.}


\index{Konashenkova T.\,D.}





{\renewcommand{\thefootnote}{\fnsymbol{footnote}}


\footnotetext[1]


{Работа выполнялась % в~рамках государственного задания <<Тео\-ре\-ти\-ко-ве\-ро\-ят\-ност\-ные и~статистические методы 


%моделирования. 0063-2019-0008>> (АААА-А19-119091990037-5) 


с~использованием инфраструктуры Центра 


коллективного пользования <<Высокопроизводительные вычисления и~большие данные>> (ЦКП 


<<Информатика>>) ФИЦ ИУ РАН (г.~Моск\-ва).}} 





\renewcommand{\thefootnote}{\arabic{footnote}}


\footnotetext[1]{Федеральный исследовательский центр <<Информатика и~управление>> Российской академии наук, 


\mbox{sinitsin@dol.ru}}


\footnotetext[2]{Федеральный исследовательский центр <<Информатика и~управление>> Российской академии наук, 


\mbox{vsinitsin@ipiran.ru}}


\footnotetext[3]{Федеральный исследовательский центр <<Информатика и~управление>> Российской академии наук, 


\mbox{ekorepanov@ipiran.ru}}


\footnotetext[4]{Федеральный исследовательский центр <<Информатика и~управление>> 


   Российской академии наук,  \mbox{tkonashenkova64@mail.ru}}











 \label{st\stat}





 \thispagestyle{headings}


 


\vspace*{-12pt}











  


  


\Abst{Построено каноническое разложение (КР) скалярного стохастического 


процесса (СтП), заданного на конечном промежутке времени, с~применением 


технологии вейв\-лет-ней\-рон\-ной сети (ВНС). Разработана архитектура 


трехслойной ВНС. Активационные функции скрытого слоя строятся на основе 


выбранного ортонормированного базиса вейвлетов с~компактными носителями. 


Для обучения ВНС применяется метод наискорейшего спуска. Дан алгоритм 


построения КР ковариационной функции заданного СтП на основе ВНС. 


Каноническому разложению ковариационной функции соответствует КР СтП 


в~виде линейной комбинации базисных вейвлетов с~коэффициентами~--- 


случайными величинами (СВ) с~нулевыми математическими ожиданиями 


и~дисперсиями, определяемыми с~помощью предлагаемого алгоритма. Дан 


иллюстративный пример. Проведено сравнение рекуррентного метода КР на 


основе вейвлетов Хаара с~предлагаемым новым методом построения КР на основе 


ВНС.} 





\KW{вейвлет; вейвлет-нейронная сеть; искусственная нейронная сеть; 


каноническое разложение; ковариационная функция; моделирование; 


стохастический процесс}


 


\DOI{10.14357\!/\!08696527240202}{YFHFIN}  





\vskip 10pt plus 9pt minus 6pt





 \thispagestyle{headings}


 


 \vspace*{-12pt}





\section{Введение}





\vspace*{-4pt}





  Как известно~[1--4], КР СтП представляет собой линейную комбинацию некоррелированных 


СВ с~коэффициентами в~виде детерминированных 


функций. Каноническое разложение СтП упрощает выполнение различных математических операций, 


особенно линейных, над случайными функциями. Это объясняется тем, что 


зависимость функции от случайного аргумента сводится к~за\-ви\-си\-мости 


функции от детерминированных, так называемых координатных, функций, 


а~следовательно, все операции над исследуемой случайной функцией сводятся к~соответствующим операциям над координатными функциями. 


Каноническое разложение СтП удобно 


для решения задач как вероятностного (аналитического), так и~статистического 


моделирования. Методы моделирования на основе КР СтП в~рамках 


корреляционной теории оказываются универсальными. При этом, следуя 


прямой и~обратной тео\-ре\-мам Пугачёва, обычно рассматривают 


моделирование КР СтП при известной ковариационной функции 


и~моделирование КР ковариационной функции при известном КР СтП. 


Создание теории КР случайной функции связано с~именами Лоэва, 


Колмогорова, Карунена и~Пугачёва в~1940--1950~гг. Развитие теории 


канонических представлений случайных функций представлено в~работах 


Пугачёва, Марченко, Синицына~[1--5]. 


  


  В последние десятилетия для анализа нестационарных СтП в~стохастических 


системах (СтС) успешно применяют вейвлеты~[6--8]. В~работе~[9] дано 


развитие теории КР СтП на основе  


вейв\-лет-тех\-но\-ло\-гий. Для построения КР нестационарного СтП были 


выбраны вейвлеты с~компактными носителями~[10].  


Вейв\-лет-ка\-но\-ни\-че\-ское разложение (ВлКР) СтП строится на основе 


коэффициентов разложения ковариационной функции СтП по 


ортонормированному двумерному базису вейвлетов с~компактными 


носителями. Вычислительные эксперименты показали, что моделирование 


нестационарного СтП с~применением ВлКР позволяет ограничиться небольшим 


числом членов КР~[11]. В~[12--16] метод ВлКР был успешно применен для 


синтеза оптимальных по различным байесовым критериям линейных 


и~нелинейных СтС. Вычислительные эксперименты подтвердили высокую 


точность и~простоту разработанных методов.


  


  В настоящее время при решении задач имитационного моделирования 


широко используются подходы, основанные на искусственных нейронных 


сетях (ИНС)~[17]. Искусственная нейронная сеть~--- это набор методов анализа данных, име\-ющих 


сетевую структуру и~об\-ла\-да\-ющих двумя характерными чертами.  


Во-пер\-вых, нейросетевой алгоритм всегда имеет распределенную сетевую 


структуру, состоящую из элементов, называемых <<формальными 


нейронами>>. Во-вто\-рых, нейросетевой алгоритм всегда представляет собой 


<<обуча\-емый алгоритм>>. Все сетевые алгоритмы уникальны. 


Отличительными чертами каждого сетевого алгоритма выступают следующие 


элементы:


  \begin{itemize}


\item число используемых формальных нейронов~--- элементов сетевой 


структуры;


\item выделение в~данной сетевой структуре упорядоченных образований, 


на\-зы\-ва\-емых слоями;


\item число использу\-емых слоев формальных нейронов;


\item наличие или отсутствие дополнительных входов и~выходов в~отдельных 


нейронных слоях;


\item наличие или отсутствие обратных связей и~внутренних итераций;


\item структура алгоритма обуче\-ния;


\item использование алгоритма обучения как в~виде отдельной стадии, так 


и~в~процессе функционирования.


\end{itemize}





  Огромное внимание уделяется построению нейросетевых алгоритмов, 


использующих вейвлеты в~качестве активационной функции скрытого 


слоя~[18, 19]. При описании таких ИНС часто используют термин 


<<вейвлон>>, который обозначает нейрон, активационной функцией которого 


служит вейвлет. Вейв\-лет-ней\-рон\-ная сеть~--- ИНС, скрытый слой 


которой состоит из вейвлонов, чьи сдвиги и~сжатия задаются при 


инициализации и~в~процессе обучения не меняются. Вейв\-лет-ней\-рон\-ные 


сети считаются относительно изученными. В~[19] показано, что использование 


многомерных вейв\-лет-функ\-ций позволяет провести более точную и~тонкую 


настройку работы сети с~отдельными входами или группами входов.


  


  В настоящей статье авторы рассматривают вопрос построения КР СтП 


в~заданном промежутке времени с~применением ВНС, учитывая 


положительный опыт построения и~применения ВлКР. В~разд.~2 описано 


построение ВлКР СтП на основе рекуррентного алгоритма~\cite{3-kon}. 


В~разд.~3 предлагается новый способ построения КР СтП на основе 


трехслойной ВНС (далее~--- КРВНС): первый слой~--- входной; второй~--- 


скрытый; третий~--- выходной. В~скрытом слое входные данные 


преобразуются в~элементы пространства признаков, заданного на основе 


ортонормированного базиса вейвлетов с~компактными носителями. 


В~выходном слое формируется выходное значение сети, которое сравнивается 


с~целевым значением. Расхождение между выходным и~целевым значениями 


используется для настройки параметров ВНС. В~разд.~4 приведен пример, 


иллюстрирующий высокую эффективность разработанного метода. Заключение 


содержит основные выводы и~направления дальнейших исследований.


  


\section{Построение вейвлет-канонического разложения стохастического процесса 


на~основе~рекуррентного алгоритма}


  


  Пусть скалярный СтП $X(t)$ для $t\hm\in T$ задан нулевым математическим 


ожиданием и~ковариационной функцией 


$$


K_X\left(t_1, t_2\right)= {\sf M}\left[ 


X(t_1)\overline{X(t_2)}\right]\enskip 


\left(t_1\in T,\ t_2\in T\right), \ 


K_X\left(t_1,t_2\right)\in \mathcal{L}^2(T\times T).


$$


%


 Построим КР заданного СтП 


$X(t)$ в~виде 


  \begin{equation}


  X(t)=\sum\limits_{v=1}^L V_v x_v(t)\,,


  \label{e1-kon}


  \end{equation}


где $V_v$~--- некоррелированные СВ с~нулевыми математическими 


ожиданиями и~дисперсиями~$G_v$; $x_v(t)$~--- координатные функции, 


удовлетворяющие соотношениям~[2, 3]:





\noindent


\begin{gather*}


V_v= \int\limits_T a_v(t) X(t)\,dt\,;\\ %\label{e2-kon}\\


x_v(t)= \fr{1}{G_v}\int\limits_T K_X(t_1,\tau) a_v(\tau)\,d\tau\,; \\


%\label{e3-kon}\\


\int\limits_T x_v(\tau) a_\mu(\tau)\,d\tau =\delta_{v\mu}.


%\label{e4-kon}


\end{gather*}


Функции $a_v(t)$ находятся из условия некоррелированности СВ~$V_v$, 


поэтому их можно определить бесчисленным множеством способов. 





Тогда, согласно прямой теореме Пугачёва, КР ковариационной функции СтП 


$X(t)$ примет вид:


\begin{equation*}


K_X(t_1,t_2) =\sum\limits^L_{v=1} G_v x_v(t_1) x_v(t_2).


%\label{e5-kon}


\end{equation*}


  


  Таким образом, для построения КР СтП $X(t)$ достаточно найти функции 


$x_v(t)$ и~$a_v(t)$ и~дисперсии~$G_v$ СВ $V_v$ ($v\hm= \overline{1,L}$).


  


  Для построения ВлКР заданного СтП $X(t)$ определим в~пространстве 


$\mathcal{L}^2(T)$ ортонормированный базис вейвлетов с~компактными 


носителями вида


  \begin{equation}


  \left\{ \varphi_{00}(t), \psi_{jk}(t)\right\},


  \label{e6-kon}


  \end{equation}


где $\varphi_{00}(t)= \varphi(t)$~--- масштабирующая функция; 


$\psi_{00}(t)\hm= \psi(t)$~--- материнский вейвлет, $\psi_{jk}(t)\hm= 


\sqrt{2^j}\,\psi(2^j t-k)$~--- вейвлеты уровня~$j$ для $j\hm= \overline{1, J}$; $k\hm= 


\overline{0,2^{j}-1}$, $J$~--- максимальный уровень вейв\-лет-раз\-ре\-ше\-ния. 


Тогда параметр $L\hm= 2^{J+1}$ равен числу базисных функций. Для удобства 


представим вейв\-лет-ба\-зис~(\ref{e6-kon}) в~виде


\begin{equation}


g_1(t)=\varphi_{00} (t)\,;\quad g_2(t)=\psi_{00}(t)\,;\quad


g_v(t)=\psi_{jk}(t),


\label{e7-kon}


\end{equation}


где


$ j=\overline{1, J}$, $k=\overline{0, 2^j-1}$, $v=2^j+k+1$.








  


  В пространстве $\mathcal{L}^2(T\times T)$ определим двумерный 


ортонормированный вейв\-лет-ба\-зис в~виде тензорного произведения двух 


вейв\-лет-ба\-зи\-сов вида~(\ref{e6-kon}) для случая, когда масштабирование по 


обеим переменным происходит одинаково. В~этом случае базисные функции 


имеют вид:


  \begin{align*}


  \Phi^A(t_1,t_2) &= \varphi_{00}(t_1) \varphi_{00}(t_2)\,;\quad 


\Psi^H(t_1,t_2)=\varphi_{00}(t_1)\psi_{00}(t_2)\,;\\


  \Psi^B(t_1,t_2) &= \psi_{00}(t_1)\varphi_{00}(t_2)\,;\quad 


\Psi^D_{jkn}(t_1,t_2)= \psi_{jk}(t_1) \psi_{jn}(t_2)\,,


  %\label{e8-kon}


  \end{align*}


где $j=\overline{1,J}$; $k,n\hm= \overline{0, 2^j-1}$. Тогда двумерное  


вейв\-лет-раз\-ло\-же\-ние заданной ковариационной функции примет вид:


\begin{multline*}


K_X(t_1,t_2) =a^x\Phi^A(t_1,t_2) +h^x\Psi^H(t_1,t_2) +b^x \Psi^B (t_1,t_2) 


+{}\\


{}+\sum\limits^J_{j=0} \sum\limits_{k=0}^{2^J-1} \sum\limits_{n=0}^{2^J-1} 


d^x_{jkn} \Psi^D_{jkn} (t_1,t_2)\,,


%\label{e9-kon}


\end{multline*}


  где 


  $$


  a^x\! =\!\int\limits_T \!\int\limits_T\! K_X(t_1,t_2) \Phi^A (t_1,t_2)\,dt_1dt_2\,;\enskip 


h^x\!=\!\int\limits_T \!\int\limits_T \!K_X(t_1,t_2) \Psi^H(t_1, t_2)\,dt_1dt_2\,;


  $$


  $$


  b^x\!=\! \int\limits_T \!\int\limits_T \!K_X(t_1,t_2) \Psi^B(t_1,t_2)\,dt_1dt_2\,;\enskip 


d^x_{jkn}\!= \!\int\limits_T \!\int\limits_T \!K_X(t_1,t_2) 


\Psi^D_{jkn}(t_1,t_2)\,dt_1dt_2\,.


  $$


  


  Введем вспомогательные параметры $k_{v\mu}$ ($v,\mu=\overline{1,L}$): 


  \begin{multline*}


  k_{11} =a^x;\enskip k_{12} =h^x; \enskip k_{21}=b^x;\enskip 


k_{22}=d_{00}^x\,;\\


  k_{v\mu}=d^x_{jkn}\enskip \left( v=2^j+k+1;\ \mu=2^j+n+1;\ 


i=\overline{1,J};\ n=\overline{0, 2^j-1}\right),


%  \label{e10-kon}


  \end{multline*}


  остальные $k_{v\mu}=0$.


  


  Дисперсии $G_v$ СВ $V_v$ вычисляются по рекуррентным формулам:


  \begin{gather*}


%  \left.


 % \begin{array}{c}


  c_{v1}=-\fr{k_{v1}}{G_1}\ \left( v=\overline{2,L}\right);\\


   c_{v\mu}=-


\fr{1}{G_\mu} \left( k_{v\mu} -\sum\limits_{\lambda=1}^{\mu-1} G_\lambda 


c_{\mu\lambda} c_{v\lambda}\right)\ \left( \mu=\overline{2,v-1};\ 


v=\overline{3,L}\right);\\


  G_1=k_{11};\enskip G_v=k_{vv}-\sum\limits_{\lambda=1}^{v-1} 


G_\lambda^x\left\vert c_{v\lambda}\right\vert^2\ \left( v=\overline{2,L}\right).


%  \end{array}


%  \right\}


%  \label{e11-kon}


  \end{gather*}


  


  Для определения координатных функций $x_v(t)$ ($v\hm= \overline{1,L}$) 


сначала вычисляются вспомогательные функции 


  \begin{gather*}


%  \left.


%  \begin{array}{c}


  h_1^x(t)=a^x\varphi_{00}(t)+b^x\psi_{00}(t);\enskip h_2^x(\overline{\tau}) =h^x 


\varphi_{00}(t)+d^x_{000}\psi_{00}(t)\,;\\


  h_v^x(t)=\displaystyle\sum\limits_{k=0}^{2^j-1} d^x_{jkn} \psi_{jk}(t)\enskip 


   \left(  v=\overline{3,L};\ v=2^j+n+1;\ n=\overline{0, 2^j-1};\ j=\overline{1,J}\right),


%  \end{array}


%  \right\}


 % \label{e12-kon}


  \end{gather*}


затем вычисляются координатные функции $x_v(t)$ последовательно по 


формулам:


\begin{gather*}


%\left.


%\begin{array}{c}


x_1(t)=\fr{1}{G_1} \,h_1^x(t);\enskip x_v(t)=\fr{1}{G_v}\left( 


\sum\limits_{\lambda=1}^{v-1} d_{v\lambda} 


h_\lambda^x(t)+h_v^x(t)\right)\,;\\


d_{v\lambda} =c_{v\lambda} +\sum\limits_{\mu=\lambda+1}^{v-1} c_{v\mu} 


d_{\mu\lambda}\ \left( \lambda=\overline{1,v-2}\right);\ d_{v,v-1} =c_{v,v-1};\ 


v=\overline{2,L}\,.


%\end{array}


%\right\}


%\label{e13-kon}


\end{gather*}


При этом функции $a_v(t)$ выражаются через базисные вейв\-лет-функ\-ции 


вида~(\ref{e7-kon}):


\begin{equation*}


a_1(t)=g_1(t),\enskip a_v(t)=\sum\limits_{\lambda=1}^{v-1} d_{v\lambda} 


g_\lambda(t) +g_v(t)\enskip \left( v=\overline{2,L}\right).


%\label{e14-kon}


\end{equation*}


  


  Разработанный алгоритм был реализован для вейвлетов Хаара в~виде 


инструментального программного обеспечения <<СтИТ-КРВЛ.1>> для среды 


MATLAB.





\section{Построение канонического разложения стохастического процесса


на~основе трехслойной вейвлет-нейронной сети}


  


  Обобщая результаты построения ВлКР СтП $X(t)$, изложенные в~разд.~2, 


будем искать его КР~(1) для координатных функций в~виде:


  \begin{equation*}


  x_v(t)=g_v(t)\enskip \left( v=\overline{1,L}\right).


  %\label{e15-kon}


  \end{equation*}


Тогда, согласно прямой теореме Пугачёва, КР ковариационной функции 


$K_X(t_1,t_2)$ примет вид: 


\begin{equation}


K_X(t_1,t_2) =\sum\limits^L_{v=1} D_v g_v(t_1) g_v(t_2).


\label{e16-kon}


\end{equation}


Тем самым задача построения КР СтП $X(t)$ сведется к~задаче аппроксимации 


ковариационной функции $K_X(t_1,t_2)$ линейной зависимостью нелинейных 


двумерных вейв\-лет-функ\-ций вида~(\ref{e16-kon}), которую можно 


вычислить с~использованием многослойного персептрона~\cite{17-kon}. После 


определения дисперсий~$D_v$ в~(\ref{e16-kon}), согласно обратной теореме 


Пугачёва, получим искомое КР СтП $X(t)$:


\begin{equation}


X(t)= \sum\limits^L_{v=1} V_v g_v(t),


\label{e17-kon}


\end{equation}


где $V_v$~--- некоррелированные СВ с~нулевыми математическими 


ожиданиями и~дисперсиями~$D_v$.


  


  Для нахождения неизвестных~$D_v$ в~(\ref{e16-kon}) задается множество 


маркированных примеров 


$$


T= \left\{ (\tau^i, d^i)\right\}^N_{i=1},\enskip 


\tau^i= \left[ t_1^i\ t_2^i\right]^{\mathrm{T}},\ d^i= K_X(t_1^i, t_2^i).


$$





 Для 


обучения многослойного персептрона с~учителем применим алгоритм 


обратного распространения ошибки~\cite{17-kon}. Он представляет собой 


градиентный метод и~обладает двумя основными свойствами: (1)~процедуру 


метода легко просчитать локально; (2)~процедура метода реализует 


градиентный спуск в~пространстве весов, при котором синаптические веса 


обновляются от примера к~примеру. 


  


  В качестве активационных функций нейронов скрытого слоя рассмотрим 


функции, построенные на основе вейв\-лет-ба\-зи\-са~(\ref{e7-kon}):


  \begin{equation*}


  \Psi_v(\tau) =\Psi_v(t_1,t_2) =g_v(t_1) g_v(t_2)\enskip \left( 


v=\overline{1,L}\right).


 % \label{e18-kon}


  \end{equation*}


Вектор признаков 


$$


\Psi(\tau) = \left[ \Psi_0(\tau)\ \Psi_1(\tau)\ \cdots\ 


\Psi_L(\tau)\right]^{\mathrm{T}}


$$


 представляет собой <<образ>>, индуцированный 


в~пространстве признаков при предъявлении вектора входного сигнала 


$\tau\hm= [t_1\ t_2]^{\mathrm{T}}$. 





  Пусть $w=[w_1\ w_2\ \cdots\ w_L]^{\mathrm{T}}$~--- вектор весовых коэффициентов, 


связывающих пространство признаков с~выходным пространством, $y\hm= 


[y_1\ y_2\ \cdots\ y_N]^{\mathrm{T}}$~--- действительный отклик системы, т.\,е.\ вектор 


оценки желаемого отклика системы $d\hm= [d_1\ d_2\ \cdots\ d_N]^{\mathrm{T}}$. При 


этом вектор весовых коэффициентов $w\hm= w(n)$ меняется при каждой 


итерации алгоритма~$n$. 


  


  Функционирование ВНС можно описать следующим образом:


  \begin{itemize}


\item входной сигнал


$$


\tau^i=\left[ t_1^i\ t_2^i\right]^{\mathrm{T}}\enskip \left(i=\overline{1,N}\right);


$$


\item скрытый слой ВНС задает пространство признаков


$$


\Psi_v(\tau^i) = 


\Psi_v(t_1^i, t_2^i)= g_v\left(t_1^i\right) g_v\left(t_2^i\right)\enskip \left(v= \overline{1,L};\ \ i= \overline{1,N}\right);


$$


\item выходной слой ВНС вычисляет отклик системы


$$


y^i(n)= 


\sum^L\limits_{v=1} w_v(n) \Psi_v(\tau^i),


$$


 или в~векторной форме 


 $$


 y^i(n)= 


\Psi^{\mathrm{T}}(\tau^i) w(n)\enskip \left(i= \overline{1,N}\right),


$$


 для итерации~$n$ алгоритма.


\end{itemize}





  Выходной сигнал $y^i(n)$ для каждого $i$-го примера сравнивается 


с~соответствующим желаемым откликом сис\-те\-мы~$d^i(n)$. Их разность 


определяет сигнал ошибки 


$$


e^i(n)= d^i - y ^i(n) = d^i - \Psi^{\mathrm{T}}(\tau^i) w(n).


$$





 Функция стоимости задается в~виде среднеквадратичной 


ошибки (СКО; Mean Squared Error~--- MSE) для всего набора примеров: 


  \begin{equation}


  E(w(n))=\fr{1}{2N} \sum\limits^N_{i=1} \left( e^i(n)\right)^2.


  \label{e19-kon}


  \end{equation}


  


  Далее надо решить задачу безусловной оптимизации функции стоимости 


$E(w)$ по отношению к~вектору весов~$w$: $E(w)\hm\to \min$. Необходимое 


условие оптимальности~--- равенство нулю вектора градиента  


ско\-рости: 


$$


\nabla E(w^*) =0\,,


$$


 где 


 \begin{multline*}


  \nabla E(w) =  \left[ \fr{\partial E}{\partial w_1}\ \fr{\partial E}{\partial w_2}\ 


\cdots\ \fr{\partial E}{\partial w_L}\right]^{\mathrm{T}}\,,


\\


  \fr{\partial E}{\partial w_l}=\fr{1}{N} \sum\limits^N_{i=1} \Psi_l(\tau^i) \left( 


  d^i-\sum\limits^L_{j=1} w_j(n)\Psi_j(\tau^i)\right)\ \left( l=\overline{1,L}\right).


 \end{multline*}


 


  Для удобства обозначим градиент $\nabla E(w) \hm= g(n)$.


  


  Методы безусловной оптимизации основаны на алгоритме 


последовательного спуска. Исходное значение $w(0)\hm=0$. Далее 


генерируется последовательность весовых векторов $w(1), w(2), \ldots$ такая, 


что значение функции стоимости уменьшается при каждой итерации 


алгоритма: $E(w(n+1))\hm< E(w(n))$. Применим \textit{метод наискорейшего 


спуска}. Корректировка весов осуществляется по формуле:


  \begin{equation}


  w(n+1)=w(n)-\eta g(n)\,,


  \label{e20-kon}


  \end{equation}


где $\eta\hm=const$~--- параметр скорости обучения, $0\hm< \eta\hm<1$. Тогда с~учетом~(\ref{e20-kon}) имеем 


$$


E(w(n+1))= E(w(n))- \eta\| g(n)\|^2.


$$


 В~общем случае шагов итерации может быть бесконечно много. Из опыта 


известно, что для выхода из процедуры адаптации весов выполняется условие 


$\vert E(w(n))\vert \hm\leq E_a$ для $E_a\hm=const$~--- параметра, 


определяющего точ\-ность приближенных вычислений.





\begin{figure}[b] %fig1


   \vspace*{-6pt}


  \begin{center}


 \mbox{%


 \epsfxsize=100.981mm 


\epsfbox{sin-1.eps}


 }


\end{center}


\vspace*{-11pt}


  \Caption{Граф передачи сигнала в~ВНС}


  \end{figure}





  В результате получаем алгоритм КР ковариационной функции на основе 


обучения ВНС методом обратного распространения ошибки. 





\pagebreak


  


  \textbf{Алгоритм.} Построение КР ковариационной функции скалярного 


СтП на основе ВНС (рис.~1).\\[-13pt]


  \begin{enumerate}[1.]


\item Задание множества обучающих примеров: вектора входного сигнала 


$\tau^i\hm= \left[ t_1^i\ t_2^i\right]^{\mathrm{T}}$ и~желаемого отклика $d^i$ ($i\hm= 


\overline{1,N}$).\\[-13pt]


\item Задание параметра скорости обучения $0\hm< \eta\hm<1$.\\[-13pt]


\item Инициализация весов. При $n=0$ $w(0)\hm=0$. Последующие шаги 


начинаются с~$n\hm= 1,2,\ldots$\\[-13pt]


\item Вычисление функций скрытого слоя:





\noindent 


$$


\Psi_v(\tau^i) = \Psi_v(t_1^i, t_2^i) = g_v(t_1^i) g_v(t_2^i)\enskip \left(v= \overline{1,L},\ \ i= 


\overline{1,N}\right).


$$





\vspace*{-2pt}





\item Вычисление фактического отклика: 





\noindent


$$


y^i(n)= \Psi^{\mathrm{T}}(\tau^i) w(n)\enskip \left(i= \overline{1,N}\right).


$$





\vspace*{-2pt}





\item Вычисление ошибки $e^i(n)\hm= d^i\hm- y^i(n)$ для каждого примера 


и~СКО для всего набора обучающих примеров~(\ref{e19-kon}). Начиная со 


второго шага, проверка условия окончания адаптации весов: 





\noindent


$$


\vert E(w(n)) 


\vert \leq E_a\,.


$$





\vspace*{-2pt}





\noindent


 Если условие выполняется, выход из алгоритма, 


в~противном случае переход к~п.~7.\\[-13pt]


\item Адаптация вектора весов. Изменяем вектор весов:





\noindent


$$


w(n+1)=w(n)-\eta g(n).


$$





\vspace*{-2pt}





\item Увеличиваем номер итерации~$n$ на~1 и~возвращаемся к~п.~5.


\end{enumerate}


  


  После обучения ВНС проводится тестирование для проверки качества 


аппроксимации.


  


  


  


  На рис.~1 представлен граф прохождения сигнала в~ВНС в~процессе 


обучения, основанного на коррекции СКО~(\ref{e19-kon}). Сигнал ошибки $e 


(n)$ инициализирует механизм управления, цель которого заключается 


в~применении последовательности корректировок к~синаптическим 


весам~$w$. Эти изменения нацелены на пошаговое приближение выходного 


сигнала~$y$ к~желаемому~$d$. На рис.~1 видно, что обучение на основе 


коррекции ошибок~--- это пример замкнутой системы с~обратной связью. Из 


теории управления известно, что устойчивость такой системы определяется 


параметром обратной связи, т.\,е.\ параметром скорости обучения~$\eta$. Для 


обеспечения устойчивости и~сходимости итеративного процесса обуче\-ния 


требуется тщательный подбор этого параметра. Пошаговая корректировка 


синаптических весов продолжается, как правило, до тех пор, пока сис\-те\-ма не 


достигнет устойчивого состояния, при котором синаптические веса 


практически стабилизируются. В~этой точке процесс обуче\-ния 


останавливается~\cite{17-kon}. При этих значениях синаптических весов 


достигается минимум функции стоимости~(\ref{e19-kon}). 


{ %\looseness=1





}


  


  Преимущества предлагаемого подхода нахождения минимума 


функции~(\ref{e19-kon}) по сравнению со стандартными методами, 


основанными на решении системы уравнений $\partial E/\partial w_l\hm= 0$ 


($l\hm= \overline{1,L}$) и~приводящими к~слож\-ным математическим 


процессам, заключаются в~простоте вычисления градиента $\nabla E(w)$ 


и~схо\-ди\-мости алгоритма при $0\hm< \eta\hm<1$. 





%\vspace*{-6pt}





\section{Пример}


  


  Рассмотрим задачу моделирования СтП $X(t)$ на промежутке $\left[ t_0, 


t_0+T\right]$ с~нулевым математическим ожиданием и~нестационарной 


ковариационной функцией $K_X(t_1,t_2) \hm= D\exp \left\{ -\alpha\vert t_2\hm-


t_1 \vert \hm+ \beta\left( t_2\hm+t_1\right)\right\}$ с~по\-мощью ВНС. Каноническое разложение СтП 


$X(t)$ строим в~виде~(\ref{e17-kon}), где СВ~$V_v$ имеют нулевые 


математические ожидания и~дисперсии $D_v\hm= w_v$. Значения 


синаптических весов~$w_v$ определяются из условия минимума СКО представления ковариационной функции $K_X(t_1,t_2)$ ее 


КР 





\noindent


  $$


  \mathrm{MSE} =E(w) =\fr{2}{N} \sum\limits_{i=1}^N \left( K_X\left(t_1^i, t_2^i\right) -


\sum\limits^L_{v=1} w_v g_v\left(t_1^i\right) g_v\left(t_1^i\right)\right)^2\,:


  $$


  $$


  \min\limits_w E(w)=\min\limits_w \left\{ \fr{2}{N} \sum\limits^N_{i=1} \left( 


K_X\left(t_1^i, t_2^i\right)-\sum\limits^L_{v=1} w_v g_v\left(t_1^i\right) g_v\left(t_2^i\right)\right)^2\right\}


  $$


для заданного набора входных данных 


$$


T=\left\{ \tau^i=\left[ t_1^i\ t_2^i\right]^{\mathrm{T}}, \ d^i=K_X(t_1^i, t_2^i);\ 


i=\overline{1,N}\right\}


$$


при следующих исходных данных: $D\hm=1$; $\alpha\hm=1$; $\beta\hm= 


0{,}001$; $t_0\hm= 0$; $T\hm=8$.





\pagebreak


  


  Для проведения вычислительных экспериментов разработано 


инструментальное программное обеспечение (ИПО) в~среде MATLAB. Вначале 


экспериментально было определено, что разработанный алгоритм обеспечивает 


максимальную скорость обучения ВНС без потери устойчивости при параметре 


скорости обучения $\eta\hm=0{,}8$ и~при числе тестовых примеров, равном 


числу базисных вейв\-лет-функ\-ций, $N\hm=L$. Далее была выявлена 


зависимость MSE от максимального уровня разрешения~$J$~базиса Хаара, 


представленная на рис.~2,\,\textit{а}. Минимальное значение, равное~0,0039, 


MSE принимает при $J\hm=2$, т.\,е.\ для 8~базисных функций Хаара. Тем 


самым было экспериментально определено оптимальное число нейронов 


скрытого слоя ВНС, рав\-ное~8. На рис.~2,\,\textit{б} представлен график 


зависимости MSE от числа~$n$ итераций алгоритма для $\eta\hm= 0{,}8$ 


и~$N\hm=L\hm=8$. MSE принимает минимальное значение, рав\-ное~0,0039, 


при $n\hm=4$.


  


  \begin{figure} %fig2


   \vspace*{1pt}


  \begin{center}


 \mbox{%


 \epsfxsize=128mm 


\epsfbox{sin-2.eps}


 }


\end{center}


\vspace*{-11pt}


  \Caption{Графики зависимости MSE аппроксимации $K_X(t_1,t_2)$ от уровня 


разрешения~$J$ базиса Хаара~(\textit{а}) и~от числа~$n$ итераций алгоритма при $J\hm=2$~(\textit{б})}


  \end{figure}


  


  На рис.~3,\,\textit{а} представлены графики точных 


значений Dtrue дисперсии $D_X(t)\hm= K_X(t,t) \hm= D\exp \{2\beta t\}$ 


и~ее оценки Dprib для выбранных значений 


параметров $\eta\hm=0{,}8$ и~$N\hm=L\hm=8$. На рис.~3,\,\textit{б} изображен 


график модуля ошибки оценки дисперсии. 


  


  


  \begin{figure} %fig3


   \vspace*{1pt}


  \begin{center}


 \mbox{%


 \epsfxsize=125.318mm 


\epsfbox{sin-3.eps}


 }


\end{center}


\vspace*{-12pt}


  \Caption{Графики точных значений Dtrue~(\textit{1}) дисперсии СтП   и~ее оценки 


Dprib~(\textit{2})~(\textit{а}) и~модуля ошибки оценки дисперсии~(\textit{б}) на отрезке $[0,8]$}


\vspace*{-9pt}


  \end{figure}


  


  В итоге КРВНС СтП примет вид:


  $$


  X(t)= \sum\limits^8_{v=1} V_v g_v(t).


  $$


  


 


    


С~по\-мощью замены переменных $\overline{t}\hm= (t\hm-t_0)/T$ выразим 


базисные функции $g_v(t)$ через вейвлеты Хаара $h_v(\overline{t})$, 


определенные на промежутке~$[0,1]$:  


$$


  g_1(t)=h_1\left(\overline{t}\right) =


  \begin{cases}


  1, & \mbox{если}\ \overline{t}\in [0,1)\,;\\


  0, & \mbox{если}\ \overline{t}\not\in [0,1)\,;


  \end{cases}


  $$


  


 


  


  


  


  \noindent


   \begin{align*}


  g_2(t)&=h_2(\overline{t}) =\begin{cases}


  1, & \mbox{если}\  \overline{t}\in \left[ 0,\fr{1}{2}\right);\\[9pt]


  -1, & \mbox{если}\ \overline{t}\in \left[ \fr{1}{2},1\right);\\


  0, &\mbox{если}\ \overline{t}\not\in [0,1]\,;


  \end{cases}


  \\


  g_3(t)&=h_3(\overline{t}) =\begin{cases}


  \sqrt{2}, & \mbox{если}\  \overline{t}\in\left[ 0,\fr{1}{4}\right);\\[9pt]


  -\sqrt{2}, &\mbox{если}\ \overline{t}\in \left[ \fr{1}{4},\fr{1}{2}\right);\\[9pt]


  0, &\mbox{если}\ \overline{t}\not\in \left[ 0,\fr{1}{2}\right);


  \end{cases}


 \\


  g_4(t) &=h_4(\overline{t}) =\begin{cases}


  \sqrt{2}, & \mbox{если}\ \overline{t}\in \left[ \fr{1}{2},\fr{3}{4}\right);\\[9pt]


  -\sqrt{2}, &\mbox{если}\ \overline{t}\in \left[ \fr{3}{4},1\right);\\[9pt]


  0, & \mbox{если}\ \overline{t}\not\in \left[ \fr{1}{2}, 1\right);


  \end{cases}


  \\


  g_i(t)&=h_i(\overline{t}) =\psi_{2k}(\overline{t}) =\begin{cases}


  2, &\mbox{если}\ \overline{t}\in \left[ \fr{k}{4},\fr{k+0{,}5}{4}\right);\\[9pt]


  -2, &\mbox{если}\ \overline{t}\in \left[ \fr{k+0{,}5}{4}\,\fr{k+1}{4}\right);\\[9pt]


  0, & \mbox{если}\  \overline{t}\notin \left[ \fr{k}{4},\fr{k+1}{4}\right)\  (k=0{,}3;\ i=4+k+1).


  \end{cases}


  \vspace*{-9pt}


  \end{align*}


  


   \begin{figure} %fig4


   \vspace*{1pt}


  \begin{center}


 \mbox{%


 \epsfxsize=128mm 


\epsfbox{sin-4.eps}


 }


\end{center}


\vspace*{-12pt}


  \Caption{Графики точных значений Dtrue~(\textit{1}) дисперсии СтП $X(t)$ и~ее оценки 


Dprib~(\textit{2})~(\textit{а}) и~модуля ошибки оценки дисперсии~(\textit{б}) при тестировании ВНС 


на отрезке~$[2,6]$}


\vspace*{-6pt}


  \end{figure}


  


  \begin{figure}[b] %fig5


 \vspace*{-6pt}


  \begin{center}


 \mbox{%


 \epsfxsize=128mm 


\epsfbox{sin-5.eps}


 }


\end{center}


\vspace*{-12pt}


\Caption{Графики точных значений Dtrue~(\textit{1}) дисперсии СтП $X(t)$ и~ее оценки 


Dprib~(\textit{2})~(\textit{а}) и~модуля ошибки оценки дисперсии~(\textit{б}) при тестировании ВНС 


на отрезке $[8,16]$}


\end{figure}


  


  Случайные величины $V_v$ ($v=\overline{1,8}$) имеют дисперсии, равные синаптическим 


весам: 





\vspace*{-2pt}





\noindent


\begin{gather*}


D[V_1] =w_1=0{,}2437;\enskip  D[V_2] = w_2= 0{,}1878;\\ 


D[V_3]=w_3= 0{,}1285;\enskip D[V_4]= w_4= 0{,}1295;\\ 


D[V_5]=w_5= 0{,}0792;\enskip D[V_6]= w_6= 0{,}0795;\\ 


D[V_7]= w_7 = 0{,}0798;\enskip D[V_8]= w_8= 0{,}0801.


\end{gather*}





\vspace*{-2pt}


  


  Обученная ВНС была протестирована для различных значений $t\hm\in 


[0,8]$, отличных от точек из набора обучения ВНС. Вычислительные 


эксперименты показали, что во всех случаях $\mathrm{MSE}\hm=0{,}0039$. На 


рис.~4,\,\textit{а} представлены графики значений Dtrue 


дисперсии $D_X(t)\hm= K_X(t,t) \hm= D\exp \{2\beta t\}$ и~ее оценки Dprib для $t\hm\in [2,6]\hm\subset [0,8]$. На 


рис.~4,\,\textit{б} изображен график модуля ошибки оценки дисперсии. 





\setcounter{figure}{6}


\begin{figure}[b] %fig7


 \vspace*{1pt}


  \begin{center}


 \mbox{%


 \epsfxsize=125.396mm 


\epsfbox{sin-7.eps}


 }


\end{center}


\vspace*{-11pt}


\Caption{Графики реализации СтП $X(t)$~(\textit{а}), точных значений Dtrue~(\textit{1}) дисперсии СтП 


$X(t)$ и~ее оценки Dprib~(\textit{2})~(\textit{б}) и~модуля ошибки оценки дисперсии~(\textit{в}) на 


отрезке $[0,8]$, полученные с~применением ИПО <<СтИТ-КРВЛ.1>>}


\vspace*{-12pt}


\end{figure}


  


  \pagebreak


  


  \


  


  \vspace*{-12pt}


  


  \setcounter{figure}{5}


  \begin{floatingfigure}{65mm} %fig6


 \vspace*{-8pt}


  \begin{center}


 \mbox{%


 \epsfxsize=59.294mm 


\epsfbox{sin-6.eps}


 }


\end{center}


\vspace*{-11pt}


\Caption{График реализации СтП  $X(t)$ на отрезке $[0,8]$, построенный на основе ВНС}


\vspace*{6pt}


\end{floatingfigure}


  


   


  Кроме того, обученная ВНС была протестирована на интервале $[8,16]$. 


В~этом случае получено $\mathrm{MSE}\hm=0{,}0040$. Результаты 


тестирования представлены на рис.~5. 


  


  





 


  На рис.~6 представлен результат моделирования СтП $X(t)$ на отрезке $[0,8]$ 


с~по\-мощью разработанной ВНС. 


  


  








  Для сравнения было построено ВлКР данного СтП $X(t)$ на отрезке $[0,8]$ 


  с~применением ИПО <<СтИТ-КРВЛ.1>>, реализующего рекуррентный алгоритм, 


изложенный в~разд.~2, для базиса Хаара с~максимальным уровнем разрешения 


$J\hm=2$. Сред\-не\-квад\-ра\-тич\-ная ошибка пред\-став\-ле\-ния ковариационной функции $K_X(t_1,t_2)$ ее КР 


равна~0,0033. На рис.~7 показаны результаты вы\-чис\-ли\-тель\-но\-го эксперимента. 


Оба метода дают одинаковый порядок точ\-ности. При этом ВНС проще 


в~математическом и~алгоритмическом плане.


  


  








\section{Заключение}





\vspace*{-4pt}


  


  Разработан новый метод построения КР скалярного СтП, заданного на конечном промежутке времени, на основе трехслойной 


ВНС.


 Скрытый слой задается на основе выбранного ортонормированного 


базиса вейв\-ле\-тов с~компактными носителями. Для обучения ВНС применяется 


метод наискорейшего спуска. 








  


  В результате вычислительных экспериментов было определено, что 


разработанный алгоритм обеспечивает максимальную ско\-рость обуче\-ния 


ВНС без потери устойчивости при параметре ско\-рости обуче\-ния,  


рав\-ном~0,8, и~при числе тестовых примеров, равном чис\-лу базисных  


вейв\-лет-функ\-ций. Сравнительный анализ двух методов построения КР: 


ВлКР (рекуррентный) и~КРВНС~--- показали одинаковую вычислительную 


точ\-ность. Преимущество КРВНС по сравнению с~ВлКР заключается 


в~представлении СтП в~виде линейной комбинации базисных  


вейв\-лет-функ\-ций со случайными коэффициентами, что упрощает 


выполнение различных математических операций. 


  


  Полученные результаты допускают обобщение на случай векторного СтП, 


а~также применение при разработке новых методов анализа и~синтеза 


СтС различного назначения. 


  


{\small\frenchspacing


{%\baselineskip=10.8pt


\addcontentsline{toc}{section}{Литература}


\begin{thebibliography}{99}


\bibitem{1-kon}


\Au{Пугачёв В.\,С.} Теория случайных функций и~ее применение к~задачам автоматического 


управления.~--- М.: Физматгиз,1962. 884~с. 


\bibitem{2-kon}


\Au{Синицын И.\,Н.} Канонические представления случайных функций и~их применение 


в~задачах компьютерной поддержки научных исследований.~--- М.: ТОРУС ПРЕСС, 2009. 


768~с.


\bibitem{3-kon}


\Au{Синицын И.\,Н.} Канонические представления случайных функций. Теория 


и~применения.~--- 2-е изд.~--- М.: ТОРУС ПРЕСС, 2023. 816~с.


\bibitem{4-kon}


\Au{Sinitsyn I.\,N.} Developing the theory of stochastic canonic expansions~/\!\!/ Pattern 


Recognition Image Analysis, 2023. Vol.~33. Iss.~4. P.~862--887.  doi: 10.1134\!/ S1054661823040429.


\bibitem{5-kon}


\Au{Марченко Б.\,Г.} Метод стохастических интегральных представлений и~его приложения 


в~радиотехнике.~--- Киев: Наукова думка, 1973. 191~с.


\bibitem{6-kon}


\Au{Meyer Y.} Analysis at Urbana. Vol.~1. Analysis in function spaces.~--- Cambridge, U.K.: 


Cambridge University Press, 1989. 436~p.


\bibitem{7-kon}


\Au{Mallat S.\,G.} Multiresolution approximations and wavelet orthonormal bases of 


$L^2(R)$~/\!\!/ T.~Am. Math. Soc., 1989. Vol.~315. Iss.~1. P.~69--87. doi: 10.2307\!/\!2001373.


\bibitem{8-kon}


\Au{Rhif M., Ben Abbes~A., Farah~I.\,R., Mart\!\'{\!$\iota$}nez~B., Sang~Y.} Wavelet transform 


application for/in non-stationary time-series analysis: A~review~/\!\!/ Appl. Sci., 2019. Vol.~9. 


Iss.~7. Art.~1345. 22~p. doi: 10.3390\!/\!app9071345. 


\bibitem{9-kon}


\Au{Синицын И.\,Н., Синицын~В.\,И., Корепанов~Э.\,Р., Конашенкова~Т.\,Д.} Дополнение 


<<Синтез нестационарных систем методом вейвлет канонических разложений>>~/\!\!/


\Au{Синицын~И.\,Н.} Канонические пред\-став\-ле\-ния случайных функций. Теория 


и~применения.~--- 2-е изд.~--- М.: ТОРУС ПРЕСС, 2023. С.~752--807.


\bibitem{10-kon}


\Au{Добеши И.} Десять лекций по вейвлетам.~--- М.--Ижевск: НИЦ <<Регулярная 


и~хаотическая динамика>>, 2004. 464~с. (\Au{Daubechies~I.} Ten lectures on wavelets.~--- 


Philadelphia, PA, USA: Society for Industrial and Applied Mathematics, 1992. 352~p.)


\bibitem{11-kon}


\Au{Синицын И.\,Н., Сергеев~И.\,В., Корепанов~Э.\,Р., Конашенкова~Т.\,Д.} 


Инструментальное программное обеспечение анализа и~синтеза стохастических систем 


высокой доступности (IV)~/\!\!/ Системы высокой доступности, 2017. Т.~13. №\,3. С.~55--69. 


EDN: ZSQHLR.


\bibitem{12-kon}


\Au{Sinitsyn I., Sinitsyn~V., Korepanov~E., Konashenkova~T.} Bayes synthesis of linear 


nonstationary stochastic systems by wavelet canonical expansions~/\!\!/ Mathematics, 2022. 


Vol.~10. Iss.~9. Art.~1517. 14~p. doi: 10.3390\!/\!math10091517.


\bibitem{13-kon}


  \Au{Синицын И.\,Н., Синицын~В.\,И., Корепанов~Э.\,Р., Конашенкова~Т.\,Д.} 


Инструментальное программное обеспечение анализа и~синтеза стохастических сис\-тем 


высокой доступности (XVII)~/\!\!/ Системы высокой доступности, 2023. Т.~19. №\,2.  


С.~5--24. doi: 10.18127\!/\!j20729472-202302-01. EDN: PCUHWT.


\bibitem{14-kon}


\Au{Sinitsyn I., Sinitsyn~V., Korepanov~E., Konashenkova~T.} Synthesis of nonlinear 


nonstationary stochastic systems by wavelet canonical expansions~/\!\!/ Mathematics, 2023. 


Vol.~11. Iss.~9. Art.~2059. 18~p.  doi: 10.3390\!/\!math11092059.


\bibitem{15-kon}


  \Au{Синицын И.\,Н., Синицын~В.\,И., Корепанов~Э.\,Р., Конашенкова~Т.\,Д.} 


Инструментальное программное обеспечение анализа и~синтеза стохастических сис\-тем 


высокой доступности (XVIII)~/\!\!/ Системы высокой доступности, 2023. Т.~19. №\,3.  


С.~18--34. doi: 10.18127\!/\!j20729472-202303-02. EDN: SQPRDO.


\bibitem{16-kon}


  \Au{Синицын И.\,Н., Синицын В.\,И., Корепанов~Э.\,Р., Конашенкова~Т.\,Д.} Синтез 


оптимальных по функциональным бейесовым критериям стохастических систем методом  


вейв\-лет-ка\-но\-ни\-че\-ских разложений~/\!\!/ Системы высокой доступности, 2023. Т.~19. 


№\,4. С.~37--50.  doi: 10.18127\!/\!j20729472-202304-03. EDN: WYQSQT.


\bibitem{17-kon}


\Au{Хайкин С.} Нейронные сети: полный курс~/ Пер. с~англ.~--- 2-е изд.~--- СПб.: 


Диалектика, 2020. 1104~с. (\Au{Haykin~S.} Neural networks. A~comprehensive foundation.~--- 


2nd ed.~--- Prentice-Hall of India Pvt. Ltd., 1999. 842~p.)


\bibitem{18-kon}


\Au{Терехов С.\,А.} Вейвлеты и~нейронные сети~/\!\!/ Научная сессия МИФИ-2001: III 


Всеросс. научн.-техн. конф. <<Нейроинформатика-2001>>: лекции по 


 нейроинформатике.~--- М.: МИФИ, 2001. С.~142--181.


\bibitem{19-kon}


\Au{Veitch D.} Wavelet neural networks and their application in the study of dynamical 


system~/\!\!/ Networks, 2005. Vol.~1. No.\,8. P.~313--320.





     \end{thebibliography}


} }











\vspace*{-6pt}





\hfill{\small\textit{Поступила в~редакцию 15.03.24}}





%\vspace*{8pt}





%\hrule





%\vspace*{2pt}








%\hrule





%\vspace*{8pt}





\newpage





\vspace*{-28pt}





\def\tit{NONSTATIONARY STOCHASTIC PROCESS MODELING BY~CANONICAL 


EXPANSION AND~WAVELET NEUTRAL NETWORK}











\def\titkol{Nonstationary stochastic process modeling by~canonical 


expansion and~WNN} %wavelet neutral network}





\def\aut{I.\,N.~Sinitsyn, V.\,I.~Sinitsyn, E.\,R.~Korepanov, and~T.\,D.~Konashenkova}





\def\autkol{I.\,N.~Sinitsyn, V.\,I.~Sinitsyn, E.\,R.~Korepanov, and~T.\,D.~Konashenkova}





\titel{\tit}{\aut}{\autkol}{\titkol}





\vspace*{-8pt}











\noindent


Federal Research Center ``Computer Science and Control'' of the Russian Academy of 


Sciences, 44-2~Vavilov Str., Moscow 119333, Russian Federation





\def\leftfootline{\small{\textbf{\thepage}


\hfill Sistemy i Sredstva Informatiki~---


Systems and Means of Informatics 2024 vol~34 no\,2}


}%


 \def\rightfootline{\small{Sistemy i Sredstva Informatiki~---


 Systems and Means of Informatics 2024 vol~34 no\,2


\hfill \textbf{\thepage}}}





\vspace*{3pt}


  


   


  


  \Abste{For scalar stochastic processes (StP) at finite time intervals and their canonical expansions (CE), 


  a~technology based on wavelet neural networks


  (WNN) is constructed. For WNN learning, the method of steepest descent was used. 


  The three-layer WNN architecture is presented. The activation functions of


  the latent layer are based on 


  chosen wavelet basis with general compact carrier. %\linebreak\vspace*{-12pt}} 


For StP covariance function, 


  a~special WNN algorithm of CE construction is developed. The covariance function CE corresponds to CE StP in the form 


  of linear combination of wavelet basis with zero mathematical expectations and variances defined by the suggested algorithm.


   A~numerical example illustrates CE WNN preference with wavelet CE.}


  


  \KWE{canonical expansion; covariance function; modeling; stochastic process; wavelet; 


wavelet neural network}


  


 





  


\DOI{10.14357\!/\!08696527240202}{YFHFIN}  





\vspace*{-10pt}








\Ack





\vspace*{-3pt}





\noindent


  The research was carried out using the infrastructure of the Shared Research Facilities ``High 


Performance Computing and Big Data'' (CKP ``Informatics'') of FRC CSC RAS (Moscow). 








%\vspace*{-3pt}





\renewcommand{\bibname}{\protect\rmfamily References}


%\renewcommand{\bibname}{\large\protect\rm References}





{\small\frenchspacing


{ %\baselineskip=12pt


\addcontentsline{toc}{section}{References}


\begin{thebibliography}{99}


\bibitem{1-kon-1}


  \Aue{Pugachev, V.\,S.} 1962. \textit{Teoriya sluchaynykh funktsiy i~ee primenenie 


k~zadacham avtomaticheskogo upravleniya} [Theory of random functions and its application to 


automatic control problems]. Moscow: Fizmatgiz. 884~p.


\bibitem{2-kon-1}


  \Aue{Sinitsyn, I.\,N.} 2009. \textit{Kanonicheskie predstavleniya sluchaynykh funktsiy i~ikh 


primenenie v~zadachakh komp'yuternoy podderzhki nauchnykh issledovaniy} [Canonical 


expansions of random functions and their applications in computer aided support]. Moscow: 


TORUS PRESS. 768~p. 


\bibitem{3-kon-1}


  \Aue{Sinitsyn, I.\,N.} 2023. \textit{Kanonicheskie predstavleniya sluchaynykh funktsiy. Teoriya 


i~primeneniya} [Canonical expansions of random functions. Theory and applications]. 2nd ed. 


Moscow: TORUS PRESS. 816~p.


\bibitem{4-kon-1}


  \Aue{Sinitsyn, I.\,N.} 2023. Developing the theory of stochastic canonic expansions. 


\textit{Pattern Recognition Image Analysis} 33(4):862--887. doi: 10.1134\!/\!S1054661823040429.


\bibitem{5-kon-1}


  \Aue{Marchenko, B.\,G.} 1973. \textit{Metod stokhasticheskikh integral'nykh predstavleniy 


i~ego prilozheniya v~radiotekhnike} [Method of stochastic integral representations and its 


applications in radio engineering]. Kiev: Naukova dumka. 191~p. 


\bibitem{6-kon-1}


  \Aue{Meyer, Y.} 1989. \textit{Analysis at Urbana. Vol.~1. Analysis in function spaces}. 


Cambridge, U.K.: Cambridge University Press. 436~p.


\bibitem{7-kon-1}


  \Aue{Mallat, S.\,G.} 1989. Multiresolution approximations and wavelet orthonormal bases of 


$L^2(R)$. \textit{T.~Am. Math. Soc.} 315(1):69--87. doi: 10.2307\!/\!2001373.


\bibitem{8-kon-1}


  \Aue{Rhif, M., A.~Ben Abbes, I.\,R.~Farah, B.~Mart\'{\!\!$\iota$}nez, and Y.~Sang.} 2019. Wavelet 


transform application for\!/\!in non-stationary time-series analysis: A~review. \textit{Appl. 


Sci.} 9(7):1345. 22~p. doi: 10.3390\!/\!app9071345. 


\bibitem{9-kon-1}


  \Aue{Sinitsyn, I.\,N., V.\,I.~Sinitsyn, E.\,R.~Korepanov, and T.\,D.~Konashenkova.} 2023. 


Dopolnenie ``Sintez nestatsionarnykh sistem metodom veyvlet kanonicheskikh raz\-lo\-zhe\-niy''] 


[Appendix ``Synthesis of nonstationary Systems by wavelet canonical 


expansions'']. Sinitsyn,~I.\,N. \textit{Kanonicheskie predstavleniya sluchaynykh funktsiy. Teoriya 


i~primeneniya} [Canonical expansions of random functions. Theory and applications]. 2nd ed. 


Moscow: TORUS PRESS. 752--807.


\bibitem{10-kon-1}


  \Aue{Daubechies, I.} 1992. \textit{Ten lectures on wavelets}. Philadelphia, PA: Society for 


Industrial and Applied Mathematics. 352~p.


\bibitem{11-kon-1}


  \Aue{Sinitsyn, I.\,N., I.\,V.~Sergeev, E.\,R.~Korepanov, and T.\,D.~Konashenkova.} 2017. 


Instrumental'noe programmnoe obespechenie analiza i~sinteza stokhasticheskikh sis\-tem 


vysokoy dostupnosti (IV) [Software tools for analysis and synthesis of stochastic systems with high 


availability (IV)]. \textit{Sistemy vysokoy dostupnosti} [Highly Available Systems] 13(3):55--69. 


EDN: ZSQHLR.


\bibitem{12-kon-1}


  \Aue{Sinitsyn, I., V.~Sinitsyn, E.~Korepanov, and T.~Konashenkova.} 2022. Bayes synthesis of 


linear nonstationary stochastic systems by wavelet canonical expansions. \textit{Mathematics} 


10(9):1517. 14~p. doi: 10.3390\!/\!math10091517.


\bibitem{13-kon-1}


  \Aue{Sinitsyn, I.\,N., V.\,I.~Sinitsyn, E.\,R.~Korepanov, and T.\,D.~Konashenkova.} 2023. 


Instrumental'noe programmnoe obespechenie analiza i~sinteza stokhasticheskikh sis\-tem 


vysokoy dostupnosti (XVII) [Software tools for analysis and synthesis of stochastic systems with 


high availability (XVII)]. \textit{Sistemy vysokoy dostupnosti} [Highly Available Systems]  


19(2):5--24. doi: 10.18127\!/\!j20729472-202302-01. EDN: PCUHWT.


\bibitem{14-kon-1}


  \Aue{Sinitsyn, I., V.~Sinitsyn, E.~Korepanov, and T.~Konashenkova.} 2023. Synthesis of 


nonlinear nonstationary stochastic systems by wavelet canonical expansions. \textit{Mathematics} 


11(9):2059. 18~p. doi: 10.3390\!/\!math11092059.


\bibitem{15-kon-1}


  \Aue{Sinitsyn, I.\,N., V.\,I.~Sinitsyn, E.\,R.~Korepanov, and T.\,D.~Konashenkova.} 2023. 


Instrumental'noe programmnoe obespechenie analiza i~sinteza stokhasticheskikh sistem 


vysokoy dostupnosti (XVIII) [Software tools for analysis and synthesis of stochastic systems with 


high availability (XVIII)]. \textit{Sistemy vysokoy dostupnosti} [Highly Available Systems] 


19(3):18--34. doi: 10.18127\!/\!j20729472-202303-02. EDN: SQPRDO.


\bibitem{16-kon-1}


  \Aue{Sinitsyn, I.\,N., V.\,I.~Sinitsyn, E.\,R.~Korepanov, and T.\,D.~Konashenkova.} 2023. 


Sintez optimal'nykh po funktsional'nym beyesovym kriteriyam stokhasticheskikh sistem metodom 


veyvlet-kanonicheskikh razlozheniy [Methodological support of functional Bayes synthesis of 


stochastic systems based on wavelet canonical expansions]. \textit{Sistemy vysokoy dostupnosti} 


[Highly Available Systems] 19(4):37--50. doi: 10.18127\!/ j20729472-202304-03. EDN: 


WYQSQT.


\bibitem{17-kon-1}


  \Aue{Haykin, S.} 1999. \textit{Neural networks. A~comprehensive foundation}. 2nd ed. 


Prentice-Hall of India Pvt. Ltd. 842~p.


\bibitem{18-kon-1}


  \Aue{Terekhov, S.\,A.} 2001. Veyvlety i~neyronnye seti [Wavelets and neural networks]. 


\textit{Nauchnaya sessiya MIFI-2001: III Vseross. nauchn.-tekhn. konf. ``Neyroinformatika-2001'': 


lektsii po neyroinformatike} [Scientific session MEPhI-2001: 3rd All-Russian Scientific and 


Technical Conference ``Neuroinformatics-2001'': Lectures on neuroinformatics]. Moscow: MIFI. 


142--181.


\bibitem{19-kon-1}


  \Aue{Veitch, D.} 2005. Wavelet neural networks and their application in the study of dynamical 


system. \textit{Networks} 1(8):313--320.


  


 \end{thebibliography}


} }





\vspace*{-6pt}





\hfill{\small\textit{Received March 15, 2024}} 





\vspace*{-10pt} 





\Contr





\vspace*{-3pt}


  


  \noindent


  \textbf{Sinitsyn Igor N.} (b.\ 1940)~--- Doctor of Science in technology, professor, Honored 


scientist of RF, principal scientist, Federal Research Center ``Computer Science and Control'' of the 


Russian Academy of Sciences, 44-2~Vavilov Str., Moscow 119133, Russian Federation; 


\mbox{sinitsin@dol.ru} 


  


  \vspace*{3pt}


 


 %\pagebreak


  


  \noindent


  \textbf{Sinitsyn Vladimir I.} (b.\ 1968)~--- Doctor of Science in physics and mathematics, 


principal scientist, Federal Research Center ``Computer Science and Control'' of the Russian 


Academy of Sciences, 44-2~Vavilov Str., Moscow 119133, Russian Federation; 


\mbox{vsinitsin@ipiran.ru}


  


  \vspace*{3pt}


  


  \noindent


  \textbf{Korepanov Eduard R.} (b.\ 1966)~--- Candidate of Science (PhD) in technology, leading 


scientist, Federal Research Center ``Computer Science and Control'' of the Russian Academy of 


Sciences, 44-2~Vavilov Str., Moscow 119133, Russian Federation; \mbox{ekorepanov@ipiran.ru}


  


  


  \vspace*{3pt}


  


  \noindent


  \textbf{Konashenkova Tatiana D.} (b.\ 1964)~--- Candidate of Science (PhD) in physics and 


mathematics, scientist, Federal Research Center ``Computer Science and Control'' of the Russian 


Academy of Sciences, 44-2~Vavilov Str., Moscow 119133, Russian Federation; 


\mbox{tkonashenkova64@mail.ru}








\label{end\stat}





\renewcommand{\bibname}{\protect\rm Литература}





  


